\subsection{Installment 5: ``Folkebladet,'' April 10, 1918}


The Congregationalists' corresponding treatment of these matters is found in the literature which they call church polity. Ditedal too is satisfied with this term: he uses it four times in his short affidavit concerning \textit{Aaben Erkl\ae{}ring} (\textit{Veiledning}, 198 ff.).

It is not I who have put this expression into their mouths. I hasten also to add that it is not I who supplied the busts on that six-foot-long constitutional shelf, of which I wrote, with the name \textit{kirkeforfatning}, or church polity. They have this name on their title-pages. It is natural enough for the average churchgoer to think that this name has chiefly to do with legal provisions, laws, and rules concerning the church. But a theologian who knows nothing more and better about the ``field'' deserves our undivided sympathy and a heartfelt wish for improvement.

3. E. and D. have also been influenced by the biblicistic tendency whose original spokesman was J.~T.~Beck. The second volume of his work Fødelære devotes about 50 pages ``The Constitution of the Free Congregation'' (the heading). In 27 pages Beck treats what he calls the ``principles'' or the ``inner constitution.'' Under the heading of the ``outer constitution'' he treats the gifts of grace. Beck's conception of the local congregation is more like that of the Norwegian reform advocates than that of the Congregationalists. But his pronounced biblicism, his strong emphasis on the gifts of grace (in this he represents an advance beyond the reform advocates and the older Congregationalists), together with his sharp judgments upon the congregations, that they had advanced no further than mission congregations, are factors that have been determinative for \textit{Levende principer}.

4. E. and D. also have much in common with the eminent Lutheran jurist Rudolph Sohm, who, more vigorously than Beck, places the emphasis on the charismatic organization (the church's organization by the gifts of grace) and upon the freedom of the Spirit in contrast to everything that is juridical. Sverdrup's words, ``with the diversity of the gifts of grace and their right use, the freedom of the Spirit wins its true triumph,'' are a genuinely Sohmian sentence.

Whereas Sverdrup places the emphasis on the use of the gifts of grace essentially in the local congregation, Sohm finds that the local congregation is altogether too narrow for this purpose. He maintains that the holy catholic church, the communion of saints, without regard to organization, is the place for their use. To be sure, Sverdrup would not have asserted that every organization which calls itself a congregation has the right to determine whether the Christian shall use his gift of grace or not (II, 314). But he held fast in principle to this, that ``without the congregation, the gifts of grace have no attachment and no purpose'' (\textit{Luth.\ Kirkeblad} 1893). He regarded, for example, ``Christian lay activity as one of the most precious gifts of grace in our time'' (there).

With respect to Sverdrup's mild conditions for the reception of members into the congregation, and his likewise mild doctrine of church discipline, this is as characteristically Sohm as it can be in one who, differing from Sohm, taught that ``the congregation is the proper form,'' and that all the glorious things which are said about the church, the communion of saints, in Sohm's brilliant, spirited work \textit{Church Law}, are to apply to the local congregation.

Sohm regards the church as wholly and entirely religious, an organism, not an organization; that is to say, the congregation is not any local organization, no organized church body, no state church, no confessional church (denomination); but it is God's people on earth, the body of Christ, the work of the Holy Spirit. ``Local congregations'' and other ecclesiastical organizations are indeed necessary, but none of them is holy. They are no more ``holy'' or more ``church'' than any other organization which men have agreed to form and maintain.

In this Sohm differs from Sverdrup, who according to \textit{Samlede Skrifter} (II, 5) called ``the holiness of the congregation'' ``the real main question in the Free Church.'' Indeed, Sverdrup went so far as to say that the cause of the American doctrinal controversy was ``that men have different conceptions of the congregation.'' ``But behind the doctrinal controversy too there has lurked a different conception of the congregation and of its composition and work'' (II, C, 7).

Sohm does not like the word church law or church polity, because he means that the church cannot be organized: the body of Christ cannot have juridical determinations imposed upon it; the outward ecclesiastical system is, strictly speaking, no church system at all, and therefore its juridical norms (all law is formal, coercive) are a contradiction of the essence of freedom. His well-known thesis concerning church law therefore runs thus: ``All church law is a contradiction of the essence of the church.'' The congregation, he teaches, is no definite empirical entity, no social concept in general; it is a spiritual reality raised above earthly norms; it withdraws itself from the dominion of ``law'' and is unfit for juridical organization. (Luther's view.)

Although Sohm rejects the expressions ``church law'' and ``church polity,'' yet he does bow to common usage and writes an epoch-making work with the title \textit{Church Law}. If one will accept Sohm's rejection of these expressions, then one must also adopt his definition of ``law'' (which jurists as a rule will not do) and understand the Free Church's main principle.

Meanwhile one cannot all at once change the name of the literature which casts the strongest light upon principles such as \textit{Levende principer}, or even upon Sohm's own ecclesiastical views. After all, it is not so much the name that matters as the benefit---the benefit of consulting literature on church polity in order to see whether \textit{Levende principer} stands or fails. For in this literature there is a wealth of contributions for and against the norms and forms that are determinative for the forms in which the church becomes manifest, its inner and outer life. Most of these contributions will be built upon Scripture or upon history or upon both. It is characteristic that the pastor J.~J.~Jansen (N\o{}sen) relates that when he wished to write about the concept of church and office, he turned to church law (he names a full dozen authors) there to ``study, understand, and arrive first at a conclusion about what the New Testament teaches concerning the office.'' It may seem to be a strange way of studying the New Testament---by a roundabout way. But Jansen doubtless knew what he was doing.

But back to Sohm. E. and D. had their ecclesiastical views essentially ready long before Sohm came forth with his in 1892. Yet it cannot be denied that Sverdrup found in Sohm a powerful support for his conception of the gifts of grace and the freedom of the Spirit. For after 1892 he was much more forceful in advancing it than before. And many of his ablest articles bear witness---linguistically as well---to influence from Sohm.

In Ditedal there are fewer traces of this influence. But Ditedal, on the other hand, had from earlier days been strongly influenced by Professor Michael Baumgarten, who---as Professor Kaunleiter points out (\textit{Herzog-Hauck's Encyclop\ae{}die} II, 461)---was a kind of forerunner of Sohm. Baumgarten could say: ``God grant we had never got a church law. We should then long ago have had a `spiritual' church order.''

Baumgarten was one of Germany's most zealous champions of free-church ideas, and believed in what he called ``the autonomy of the saints.'' He felt nothing but aversion toward the connection between state and church, laid stress upon the freedom and inwardness of faith, hated all coercion and legalism. He taught that the church ought to use the apostolic writings as its standard also for church government: there ought to be an apostolic polity. It was chiefly the latter that caused the church historian Kliefoth to secure from the Duke of Mecklenburg Baumgarten's dismissal as theological professor at the University of Rostock. In his understanding of the Acts of the Apostles Ditedal was markedly influenced by Baumgarten's work on apostolic history (1858), which is not difficult to demonstrate.

These four currents---the reform movement, Congregationalism, the Beckian, and the Sohmian---the authors of the \textit{Principles} have taken into account. The kinship is greatest, however, with the Congregationalists, and next with Sohm, despite the fact that Sohm denies principle no.~1.

1.

I do not hesitate to say that I cannot see how anyone can rightly assess Sverdrup's ecclesiastical views unless he acquaints himself with the literature on church polity of the older Congregationalists and with Sohm's works on church law.

To suggest something of this sort---that one must read certain writings, not authored by E. and D., in order rightly to understand these writings authored by E. and D.---can hardly do otherwise than have an irritating effect. I know that very well. But as one who has ``made the attempt,'' and thus speaks from a little experience, I am free enough to say it. For among us too there are those who read the principles and interpret them as a Jew with the books of Moses reads his Old Testament, absolutely certain that they have the right understanding, that they see the guiding thread running through it, and that they have no need of guidance.

I especially single out Sohm. He can indeed offer those who are not too proud for it a helping hand. If the recommendation I give him does not suffice, then perhaps the following from a pastor in Norway will be acceptable:

\begin{quote}
``It is \ldots{} Sohm's great merit to have vindicated the great significance of the pneumatic-charismatic element (\textit{pneuma} = spirit, \textit{charisma} = gift of grace) in the constitution of the ancient church. He has found the right point of departure for understanding what is characteristic and unique in the church's formation as a community, and thereby has become able, with unsurpassed clarity and sharpness, to set forth the constitutional-historical problems which here call for the attention of research.

``One's heart can be seized by the lofty idealism in Sohm's thoughts. One feels how he suffers under the church's dualism, how it pains him to see the element of law bring a foreign and dangerous spirit into the church's development, how he longs for harmony and clarity. And in all his work one particularly notices how he burns with zeal to bring God glory and let his Spirit be the church's only Lord and ruler. He would cleanse the alien elements out of God's house and make earnest of the fact that this shall be the Lord's body and a spiritual temple built of living stones.''
\end{quote}

How strongly this recalls Sverdrup's statement: ``That which is really the principal matter in the free congregation, and which we cannot stress often enough, is the unfolding and right use of the gifts of grace.'' Or, ``with the diversity of the gifts of grace and their right use, the freedom of the Spirit wins its true triumph'' (S.~S.~II).

The longer Sverdrup wrote, the more he tended toward Sohm. And it is possible that Sohm's new work of 1914, which represents an intense labor continued all the way from 1892, when his first volume on church law appeared, might have led Sverdrup to modify his view with respect to principle no.~1. The line of thought of the two men is very similar: agreement as to the consequences seems unavoidable.

One of Sohm's colleagues defines the concept of the church as follows:

\begin{quote}
``\S 1. The church is a religious reality, God's people, the body of Christ. It has its citizenship in heaven and has there its home. But in this church God's redeeming will is carried through in mankind here on earth. \S 2. Therefore the church here on earth appears in the form of a community: a communal life. This form, however imperfect it may be, is essential for the church. Consequently the communal aspect is more than a `support and help'; it belongs to the very concept of the church, and is the form in which the church of faith is realized, in so far as it can be realized here in the world.''
\end{quote}

Sohm and Sverdrup would agree in \S 1.

Sohm would reject \S 2 and say that this ``form'' is not essential for the church, does not belong to the very concept of the church, but is only a support and help. Sverdrup would have accepted \S 2, but on the condition that ``form'' means the local congregation, and that all other outward ecclesiastical organizations are ``support and help,'' or barrier and hindrance, according to circumstances.

